\begin{center}
    \Large{\textbf{Devoir maison - piles et files} }
\end{center}
On s'intéresse à la fonction tri\_pair qui prend en paramètre une pile de nombre entier et qui renvoie une pile ne contenant que les nombres pairs de cette même pile.

\begin{lstlisting}[language=Python]
def tri_pair(p):
    p_pair = []
    while p != []:
        v = p.pop()
        if v%2 == 0 :
            p_pair.append(v)
    return p_pair
\end{lstlisting}

\begin{Exercise}[title={Pile de nombres pairs}]

\Question Quelle politique d'accès suivent les piles ?\\
%\emptybox{1cm}

\Question Que teste la condition de la ligne 3 du code ?\\
%\emptybox{1cm}

\Question Compléter le tableau d'avancement des variables pour un appel de la fonction tri\_pair([5, 2, 6, 7, 1]).\\
%\emptybox{1cm}

\begin{center}
\large
\begin{tabular}{c|c|c|c|c|}
\textbf{Tour} & v & p & v\%2==0  & p\_pair  \\
\hline
0 &  &  &  & \\
1 &  &  &  & \\
2 &  &  &  & \\
3 &  &  &  & \\
4 &  &  &  & \\
\end{tabular}
\end{center}

\Question Dessiner les piles p et p\_pair à chaque fin d'exécution 
\begin{itemize}
    \item tri\_pair([-1, 144, -10, 23])
    \item tri\_pair([233, 677, 99, 3, 73, 75, 91])
    \item tri\_pair([12, 5, 8, 6, 2])
\end{itemize}
%\emptybox{5cm}\\

\Question Qu'observez-vous sur l'ordre des piles retournées ? Que faudrait-il faire pour régler le problème ? (ne pas le coder)\\
%\emptybox{2cm}


\Question On souhaite maintenant modifier la fonction, pour qu'elle renvoie une pile contenant les nombres pairs et une pile contenant les nombres impairs. Compléter les trous ci-dessous pour suivre ce nouveaux comportements.\\
\begin{lstlisting}[language=Python]
def tri_pair(p):
    p_pair = []
    _ _ _ _ _ _  _ _ _ 
    while p != []:
        v = p.pop()
        if v%2 == 0 :
            p_pair.append(v)
        _ _ _ _ _ _ _  _
            _ _ _ _ _ _ _ _ _ 
    return _ _ _ _ _ _ 
\end{lstlisting}


\end{Exercise}